\section{Discussion}

\subsection{Key Findings}

\textbf{Finding 1: Benchmark score trajectories are logistically distributed---decisively, not marginally.}
$\Delta$AIC values of $-194$ to $-250$ versus linear (and $-113$ to $-357$ versus power-law) are 11--25$\times$ the decisive evidence threshold.
The linear model's $R^2$ of $\sim$0.53--0.67 confirms it is a genuinely poor description of leaderboard dynamics.
This has an immediate implication: ``our method achieves X\% above previous state-of-the-art'' means something fundamentally different depending on whether the benchmark is in rapid growth or plateau saturation.
Our pipeline provides infrastructure to make this distinction quantitative.

\textbf{Finding 2: Benchmark-specific overfitting inflected before benchmarks launched.}
Negative $t_0$ values for both GLUE and SuperGLUE indicate the inflection point occurred before leaderboard tracking began.
The most compelling explanation: BERT, XLNet, and related models published in late 2018 had already captured essentially all benchmark-exploitable NLU signal.
The ``rapid NLP progress'' narrative of 2018--2021 may be better described as efficient exploitation of a capability ceiling that pretraining had already established---not genuine rapid improvement.
This finding requires independent verification with verified pre-launch submission timestamps.

\textbf{Finding 3: The saturation detector is precise retrospectively but fails prospectively.}
3-month and 5-month retrospective accuracies demonstrate correct identification of community-recognized saturation moments.
The P3 failure clarifies scope: this is a retrospective monitoring tool.
The path to early warning requires shorter lookback (1--3 months) or benchmarks with faster saturation dynamics.

\subsection{Limitations}

\textbf{L1 (Data infrastructure)}: Papers With Code API deprecated; curated fallback used.
The pipeline's automation claim requires HuggingFace API adaptation.
Methodology validity is unaffected.
\textbf{L2 (Two benchmarks)}: Generalization to vision, multilingual, or generation benchmarks requires further empirical work.
\textbf{L3 (Synthetic boundary testing)}: H-C1 finding based on synthetic benchmarks; real small-benchmark validation needed.
\textbf{L4 (Self-reporting bias)}: Fitted $K$ values ($\sim$0.89) consistent with human parity, suggesting bias is small; not directly tested.
\textbf{L5 (Negative $t_0$ interpretation)}: Mechanistic interpretation plausible but unverified; competing explanation (logistic backward extrapolation) cannot be ruled out without verified pre-launch data.

\textbf{L6 (Ground truth circularity)}: The saturation dates used for validation (Sept 2019 for GLUE from the SuperGLUE paper \citep{wang2019superglue}; Jun 2021 for SuperGLUE from BIG-bench \citep{srivastava2022bigbench}) are community-documented events recorded retroactively in published papers, not independently measured saturation timestamps.
There is an inherent circularity: we validate a detector designed to systematize community judgment by comparing against that same community judgment.
An ideal validation would employ blind annotation of saturation dates by annotators who had not observed community discussions, or prospective prediction prior to community action.
We report this as a structural limitation of any validation approach for this problem: the phenomenon we aim to detect was not independently instrumented at the time it occurred.

\subsection{Broader Impact}

Automated saturation detection enables proactive benchmark retirement, contextualized paper review, and historical saturation auditing.
The primary risk of misuse is premature benchmark retirement; we recommend saturation scores be used alongside behavioral evaluation, out-of-distribution testing, and downstream task performance rather than as a sole criterion.
